\documentclass[11pt,a4paper]{article}
\usepackage[T1]{fontenc}
\usepackage{lmodern}
\usepackage[margin=27mm]{geometry}
\usepackage{amsmath,amssymb,amsthm,mathtools}
\usepackage{microtype}
\usepackage{booktabs,array}
\usepackage{enumitem}
\usepackage[hidelinks]{hyperref}
\setlength{\parskip}{5pt plus 1pt}
\setlength{\parindent}{0pt}
\setlength{\emergencystretch}{3em}
\setlist{itemsep=2pt,topsep=3pt}
\raggedbottom
\newtheorem{theorem}{Theorem}[section]
\newtheorem{lemma}[theorem]{Lemma}
\newtheorem{corollary}[theorem]{Corollary}
\newtheorem{proposition}[theorem]{Proposition}
\theoremstyle{definition}
\newtheorem{definition}[theorem]{Definition}
\newtheorem{example}[theorem]{Example}
\theoremstyle{remark}
\newtheorem{remark}[theorem]{Remark}
\newcommand{\Z}{\mathbb{Z}}
\newcommand{\N}{\mathbb{N}}
\newcommand{\Q}{\mathbb{Q}}
\newcommand{\F}{\mathbb{F}}
\newcommand{\ps}[1]{#1[[x]]}
\newcommand{\coeff}[2]{[x^{#1}]\,#2}
\newcommand{\hl}[1]{\textbf{#1}}
\newcommand{\para}[1]{\par\medskip\noindent\textbf{#1}\enspace}
\hypersetup{pdftitle={The Catalan numbers inside the equation A(x - x²) = x + A(x)²: a proof of a conjecture in OEIS A393172},pdfauthor={Dev D. Goyal},pdfsubject={OEIS A393172}}
\title{The Catalan numbers inside the equation $A(x-x^2)=x+A(x)^2$\\[6pt]\large A proof of a conjecture in OEIS A393172}
\author{Dev D. Goyal}
\date{October 2026}
\begin{document}
\maketitle
\begin{abstract}
There is exactly one power series $A(x)=x+2x^2+8x^3+42x^4+\cdots$ that satisfies the equation $A(x-x^2)=x+A(x)^2$.
Its coefficients $a(n)$ form the sequence A393172 of the On-Line Encyclopedia of Integer Sequences (OEIS).
Paul D.~Hanna conjectured that, for $n>1$, the coefficient $a(n)$ leaves remainder $2$ on division by $4$ when $n$ is a
power of $2$, and is divisible by $4$ otherwise. We prove this conjecture. The key observation is that, once every
coefficient is replaced by its remainder modulo $4$, the series $A(x)$ becomes $2C(x)-x$, where
$C(x)=x+x^2+2x^3+5x^4+\cdots$ is the generating function of the Catalan numbers. The conjecture then follows from the
classical fact that the coefficient of $x^n$ in $C(x)$ is odd exactly when $n$ is a power of $2$. The result has also been checked
by computer and proved formally in the Lean proof assistant.
\end{abstract}
\begin{center}\begin{minipage}{0.88\linewidth}\small
\textbf{How this paper was produced.} This paper reports the result of an AI research effort directed by
Dev~D.~Goyal. ChatGPT (OpenAI), using the GPT-6 model with Pro reasoning and working autonomously from prompts by
Dev~D.~Goyal, found the proof, wrote its Lean~4 formalization, and wrote a first draft of the paper.
Dev~D.~Goyal chose the problems and asked for proofs, formal verification and revisions, but did not supply proof
ideas. Claude Opus~5.5, used with Claude Code, then checked the work independently: it rebuilt the formal proofs from
source and confirmed that they rely only on Lean's three standard axioms, recomputed the sequence from its definition
compared it with the published terms and tested every claim numerically, and read every step of the proofs. It also wrote the present
version of this paper.
Dev~D.~Goyal has reviewed the proofs and is responsible for them; this review is not an independent expert validation.
The paper has not been peer reviewed.
The section ``About this paper'' at the end gives the details.
\end{minipage}\end{center}

\section{Introduction}

\subsection{The sequence}\label{sec:sequence}
The \emph{On-Line Encyclopedia of Integer Sequences} (OEIS)~\cite{oeisdb} is a large public database of integer
sequences. Many
entries also record \emph{conjectures}: patterns that the contributor observed in the data but did not prove.

This paper is about entry A393172~\cite{oeis}, contributed by Paul~D.~Hanna in February 2026. The entry reads, in part:
\begin{quote}\small\ttfamily\raggedright G.f. A(x) satisfies A(x-x\^{}2) = x + A(x)\^{}2.\par\smallskip Conjecture: for n > 1, a(n) == 2 (mod 4) iff n is a power of 2, otherwise a(n) is divisible by 4.\par\smallskip (3) A(x) = C(x) + A(C(x))\^{}2, where C(x) = x + C(x)\^{}2 is a g.f. of the Catalan numbers (A000108).\par\smallskip\hfill\textrm{--- OEIS A393172, Paul D. Hanna, February 2026}\end{quote}

Here ``G.f.''\ stands for \emph{generating function}: the terms $a(1),a(2),a(3),\dots$ of the sequence are the
coefficients of the power series $A(x)=a(1)\,x+a(2)\,x^2+a(3)\,x^3+\cdots$. To avoid confusion with the equation numbers
of this paper, we refer to the formula numbered (3) in the entry as (E3). In the equation, $A(x-x^2)$ is the series
obtained by substituting the polynomial $x-x^2$ for $x$ in $A$, and $A(x)^2$ is the square of $A(x)$. Section~\ref{sec:background}
explains these operations. The equation determines all the coefficients, one at a time (Theorem~\ref{thm:main}). The
first terms, starting at $n=1$, are 1, 2, 8, 42, 260, 1808, 13744, 112290, 974880, 8922276, \dots

Formula (E3) of the entry brings in the \emph{Catalan numbers} $1,1,2,5,14,42,132,\dots$ (OEIS A000108), one of the most
common sequences in combinatorics. For example, they count the full binary trees with $1,2,3,4,5,\dots$ leaves, that
is, the rooted trees in which every internal vertex has exactly two children, a left one and a right one. Their
generating function is
\[
C(x)=x+x^2+2x^3+5x^4+14x^5+42x^6+\cdots,
\]
and we write $c_n$ for its coefficient of $x^n$. This $C$ is the only series without constant term that satisfies
$C=x+C^2$. (In the language of trees: a tree is either a single leaf, counted by $x$, or a root with two smaller trees
attached, counted by $C^2$.)

\subsection{The question}
Hanna's conjecture says that, for $n>1$, the remainder of $a(n)$ on division by $4$ is $2$ when $n$ is a power of $2$,
and $0$ otherwise. The first row of the table shows the computed remainders of $a(n)$. The indices printed in bold are the powers of $2$.
The second row anticipates the main result: it shows the remainders modulo $4$ of the coefficients of $2C(x)-x$, which
are $1$ for $n=1$ and $2c_n$ for $n>1$.
\begin{center}\small\setlength{\tabcolsep}{3.2pt}
\begin{tabular}{>{\raggedright\arraybackslash}p{8.4em}>{\centering\arraybackslash}p{1.15em}>{\centering\arraybackslash}p{1.15em}>{\centering\arraybackslash}p{1.15em}>{\centering\arraybackslash}p{1.15em}>{\centering\arraybackslash}p{1.15em}>{\centering\arraybackslash}p{1.15em}>{\centering\arraybackslash}p{1.15em}>{\centering\arraybackslash}p{1.15em}>{\centering\arraybackslash}p{1.15em}>{\centering\arraybackslash}p{1.15em}>{\centering\arraybackslash}p{1.15em}>{\centering\arraybackslash}p{1.15em}>{\centering\arraybackslash}p{1.15em}>{\centering\arraybackslash}p{1.15em}>{\centering\arraybackslash}p{1.15em}>{\centering\arraybackslash}p{1.15em}}
\toprule
$n$ & \hl{1} & \hl{2} & 3 & \hl{4} & 5 & 6 & 7 & \hl{8} & 9 & 10 & 11 & 12 & 13 & 14 & 15 & \hl{16} \\ \midrule
$a(n)\bmod 4$ & \hl{1} & \hl{2} & 0 & \hl{2} & 0 & 0 & 0 & \hl{2} & 0 & 0 & 0 & 0 & 0 & 0 & 0 & \hl{2} \\
$\coeff{n}{(2C{-}x)}\bmod 4$ & \hl{1} & \hl{2} & 0 & \hl{2} & 0 & 0 & 0 & \hl{2} & 0 & 0 & 0 & 0 & 0 & 0 & 0 & \hl{2} \\
\bottomrule
\end{tabular}
\\[4pt]
\begin{tabular}{>{\raggedright\arraybackslash}p{8.4em}>{\centering\arraybackslash}p{1.15em}>{\centering\arraybackslash}p{1.15em}>{\centering\arraybackslash}p{1.15em}>{\centering\arraybackslash}p{1.15em}>{\centering\arraybackslash}p{1.15em}>{\centering\arraybackslash}p{1.15em}>{\centering\arraybackslash}p{1.15em}>{\centering\arraybackslash}p{1.15em}>{\centering\arraybackslash}p{1.15em}>{\centering\arraybackslash}p{1.15em}>{\centering\arraybackslash}p{1.15em}>{\centering\arraybackslash}p{1.15em}>{\centering\arraybackslash}p{1.15em}>{\centering\arraybackslash}p{1.15em}>{\centering\arraybackslash}p{1.15em}>{\centering\arraybackslash}p{1.15em}}
\toprule
$n$ & 17 & 18 & 19 & 20 & 21 & 22 & 23 & 24 & 25 & 26 & 27 & 28 & 29 & 30 & 31 & \hl{32} \\ \midrule
$a(n)\bmod 4$ & 0 & 0 & 0 & 0 & 0 & 0 & 0 & 0 & 0 & 0 & 0 & 0 & 0 & 0 & 0 & \hl{2} \\
$\coeff{n}{(2C{-}x)}\bmod 4$ & 0 & 0 & 0 & 0 & 0 & 0 & 0 & 0 & 0 & 0 & 0 & 0 & 0 & 0 & 0 & \hl{2} \\
\bottomrule
\end{tabular}
\\[4pt]
\begin{tabular}{>{\raggedright\arraybackslash}p{8.4em}>{\centering\arraybackslash}p{1.15em}>{\centering\arraybackslash}p{1.15em}>{\centering\arraybackslash}p{1.15em}>{\centering\arraybackslash}p{1.15em}>{\centering\arraybackslash}p{1.15em}>{\centering\arraybackslash}p{1.15em}>{\centering\arraybackslash}p{1.15em}>{\centering\arraybackslash}p{1.15em}>{\centering\arraybackslash}p{1.15em}>{\centering\arraybackslash}p{1.15em}>{\centering\arraybackslash}p{1.15em}>{\centering\arraybackslash}p{1.15em}>{\centering\arraybackslash}p{1.15em}>{\centering\arraybackslash}p{1.15em}>{\centering\arraybackslash}p{1.15em}>{\centering\arraybackslash}p{1.15em}}
\toprule
$n$ & 33 & 34 & 35 & 36 & 37 & 38 & 39 & 40 & 41 & 42 & 43 & 44 & 45 & 46 & 47 & 48 \\ \midrule
$a(n)\bmod 4$ & 0 & 0 & 0 & 0 & 0 & 0 & 0 & 0 & 0 & 0 & 0 & 0 & 0 & 0 & 0 & 0 \\
$\coeff{n}{(2C{-}x)}\bmod 4$ & 0 & 0 & 0 & 0 & 0 & 0 & 0 & 0 & 0 & 0 & 0 & 0 & 0 & 0 & 0 & 0 \\
\bottomrule
\end{tabular}
\end{center}

\subsection{What this paper proves}
\textbf{Main result} (Theorem~\ref{thm:main}). \emph{Modulo $4$, the series $A(x)$ equals $2C(x)-x$. In other words,
$a(n)-2c_n$ is divisible by $4$ for every $n>1$.}

Since $2c_n$ leaves remainder $2$ on division by $4$ when $c_n$ is odd, and remainder $0$ when $c_n$ is even, the
remainder of $a(n)$ depends only on whether $c_n$ is odd. A classical fact (Lemma~\ref{lem:catalan}) says that $c_n$ is
odd exactly when $n$ is a power of $2$. Together these prove Hanna's conjecture (Corollary~\ref{cor:conj}).

\subsection{What was known}
The OEIS entry records several equivalent forms of the defining equation, among them formula (E3) quoted above, and a
recurrence for $a(n)$ due to Seiichi~Manyama. The parity of the Catalan numbers is classical; see for
example~\cite{deutschsagan}. We are not aware of an earlier proof of the conjecture. The companion entry
A393171~\cite{oeis171}, also by Hanna, contains a conjecture modulo $8$ about the series $x+A(x)^2$; it is proved in a
separate paper of this collection, using the result proved here.

\subsection{Motivation}
Equations in which an unknown power series appears in a substitution, either as the inner or as the outer series,
are called \emph{functional equations}. They arise naturally when counting objects that are built recursively, such as trees,
and the OEIS contains many sequences defined in this way. Their coefficients usually grow quickly and look irregular.
Their remainders modulo small powers of $2$, however, often show patterns, and explaining such a pattern is a
way of understanding the sequence. In this case the explanation is that the Catalan numbers are hidden inside $A$: the
substitution $x\mapsto x-x^2$ is exactly the one that the Catalan series undoes. The proof finds a series that solves the
equation modulo $4$ and then uses the uniqueness of the solution; the same method is used in other papers of this
collection.

\subsection{How to read this paper}
Section~\ref{sec:background} explains the notions we need: power series, substitution, remainders, and the
Catalan numbers modulo $2$. Section~\ref{sec:results} states the results precisely, and Section~\ref{sec:proofs} proves
them. Section~\ref{sec:checks} describes how the results were checked, by computer and by a formal proof in Lean.
Section~\ref{sec:discussion} discusses the scope of the result and some open questions.
\section{Background}\label{sec:background}
This section collects the definitions and facts used later. Readers who know formal power series and congruences can skip it and return when needed.
\subsection{Formal power series}
A \emph{formal power series} is an expression
\[
U(x)=u_0+u_1x+u_2x^2+u_3x^3+\cdots ,
\]
determined by its sequence of coefficients $u_0,u_1,u_2,\dots$; see~\cite{niven,stanley}. No number is ever
substituted for $x$, so questions of convergence do not arise. We write $\coeff{n}{U}=u_n$ for the coefficient of $x^n$ in $U$; two series are equal
when all their coefficients are equal.

Series are added coefficient by coefficient and multiplied like polynomials:
\[
\coeff{n}{(UV)}=u_0v_n+u_1v_{n-1}+\cdots+u_nv_0 .
\]
Each coefficient of a product is a finite sum. For example, $(1-x)(1+x+x^2+x^3+\cdots)=1$, so $1+x+x^2+\cdots$ is
the \emph{multiplicative inverse} of $1-x$, written $\frac{1}{1-x}$. More generally, every series $D$ with constant term
$1$ has a unique multiplicative inverse $\frac1D$, whose coefficients are determined one at a time by $D\cdot\frac1D=1$
using only addition, subtraction and multiplication. Whenever a quotient of power series appears in this paper, its
denominator has constant term $1$, so the quotient is again a power series.

A \emph{commutative ring} is a set with an addition, a subtraction and a multiplication that satisfy the usual rules
of arithmetic; examples are $\Z$, $\Q$ and the integers modulo $m$. For a commutative ring $R$ we write $\ps{R}$ for
the set of power series with coefficients in $R$; for instance, writing $\F_2$ for the integers modulo $2$, $\ps{\F_2}$
consists of the series with coefficients modulo $2$. We write
$U(0)$ for the constant term $u_0$ of $U$, and $U=O(x^d)$ when the coefficients of $1,x,\dots,x^{d-1}$ in $U$ are all
zero. Thus $U-V=O(x^d)$ means that $U$ and $V$ have the same coefficient of $x^n$ for every $n<d$. A series with
constant term $0$ and coefficient of $x$ equal to $1$, that is, a series of the form $x+u_2x^2+u_3x^3+\cdots$, is
called \emph{normalized}.

\subsection{Substituting one series into another}
If $V(x)$ has constant term $0$, we can substitute $V$ for $x$ in $U$:
\[
U(V(x))=u_0+u_1V(x)+u_2V(x)^2+u_3V(x)^3+\cdots .
\]
This is called \emph{composition}, also written $U\circ V$. It makes sense because $V(x)^k$ has no terms below $x^k$,
so only the finitely many terms with $k\le n$ affect the coefficient of $x^n$. For example, if $U=x+x^2$ and $V=x^2$, then
$U(V(x))=x^2+x^4$. Composition is associative: $U\circ(V\circ W)=(U\circ V)\circ W$.

Substituting a series into itself repeatedly gives its \emph{iterates}: $U^{\circ2}(x)=U(U(x))$,
$U^{\circ3}(x)=U(U(U(x)))$, and so on, with $U^{\circ0}(x)=x$ and $U^{\circ1}=U$. They must not be confused with
\emph{powers} $U^2=U\cdot U$, $U^3=U\cdot U\cdot U$. (The OEIS sometimes writes $A^k(x)$ for the $k$-th iterate. In
this paper the small circle always means an iterate, and a plain exponent always means a power.)

\subsection{Undoing a substitution}
A normalized series can always be ``undone'' by another normalized series.

\begin{lemma}[Compositional inverse]\label{lem:inverse}
For every normalized series $U$ there is exactly one normalized series $V$ with $U(V(x))=x$, and this $V$ also
satisfies $V(U(x))=x$. If $U$ has integer coefficients, so does $V$.
\end{lemma}

We call $V$ the \emph{compositional inverse} of $U$. The OEIS calls it the \emph{series reversion} of $U$. For example,
the compositional inverse of $x-x^2$ is $C(x)=x+x^2+2x^3+5x^4+14x^5+\cdots$, whose coefficients are the Catalan numbers.
Indeed $C-C^2=x$, which says exactly that substituting $C$ into $x-x^2$ gives $x$.

\begin{proof}
Write $V=x+v_2x^2+v_3x^3+\cdots$ with unknown coefficients. In
$U(V)=V+u_2V^2+u_3V^3+\cdots$, the coefficient of $x^n$ (for $n\ge2$) equals $v_n$ plus terms that involve only
$v_2,\dots,v_{n-1}$. Setting it equal to $0$ therefore determines $v_n$ from the earlier coefficients, using only
additions, subtractions and multiplications. So there is exactly one $V$ with $U\circ V=x$, and it has integer
coefficients if $U$ does. Applying the same construction to $V$ gives a normalized $W$ with $V\circ W=x$. Then
$U=U\circ(V\circ W)=(U\circ V)\circ W=W$, so $V\circ U=V\circ W=x$ as well.
\end{proof}

\subsection{Remainders and congruences}
For an integer $m\ge2$ we write $a\equiv b\pmod m$ (``$a$ is congruent to $b$ modulo $m$'') when $a-b$ is divisible by
$m$, that is, when $a$ and $b$ leave the same remainder after division by $m$. The remainder of $a$ on division by
$m$ is called its \emph{residue} modulo $m$; it is one of $0,1,\dots,m-1$.

Addition and multiplication are compatible with congruences, so the residues $0,1,\dots,m-1$, added and multiplied
modulo $m$, form a ring, written $\Z/m\Z$. When $m=p$ is a prime we also write $\F_p$. In $\F_2=\{0,1\}$ we have $1+1=0$,
so addition and subtraction coincide.

For power series with integer coefficients, $U\equiv V\pmod m$ means that \emph{every} coefficient of $U-V$ is
divisible by $m$. Equivalently, the two series become equal when every coefficient is replaced by its residue; we call
this \emph{reducing} the series modulo $m$. The key fact is the following.

\emph{Reduction modulo $m$ is compatible with all the operations above.} To reduce a sum, a product, a substitution, or
an inverse, we may reduce the ingredients first and then compute in $\Z/m\Z$. This holds because every
coefficient of the result is obtained from finitely many coefficients of the ingredients by adding, subtracting and
multiplying, and these operations are compatible with remainders.

\para{Rings.} Some statements below hold over every commutative ring, for example $\Z$, $\Q$ and $\Z/m\Z$. No
division is used in them; this matters because in $\Z/m\Z$ division is not always possible (for instance, $2$ has no inverse in
$\Z/4\Z$).

\subsection{The method}
The main proofs follow three steps.
\begin{enumerate}
\item \emph{Uniqueness.} Show that the defining equation determines the coefficients one at a time, using only
addition, subtraction and multiplication. Then the equation has exactly one solution of the required shape, in every
commutative ring, including $\Z/4\Z$.
\item \emph{Reduction.} Reducing the integer solution modulo $4$ gives a solution of the reduced equation,
because reduction is compatible with all operations.
\item \emph{A candidate solution.} Find an explicit series that satisfies the reduced equation. By uniqueness, it is
the reduction of the integer solution, and its coefficients can be read off.
\end{enumerate}
Finding the candidate in step~3 is usually the main difficulty.
\subsection{The Frobenius identity}
\begin{lemma}[Frobenius identity]\label{lem:frobenius}
Let $p$ be a prime. For every power series $U$ with coefficients in $\F_p$, we have $U(x)^p=U(x^p)$.
\end{lemma}

\begin{proof}
In $\F_p$, the binomial theorem simplifies to $(a+b)^p=a^p+b^p$, because every binomial coefficient
$\binom pi$ with $0<i<p$ is divisible by $p$. Also $c^p=c$ for every $c\in\F_p$ (Fermat's little theorem). So
$\bigl(\sum_nu_nx^n\bigr)^p=\sum_nu_n^px^{pn}=\sum_nu_nx^{pn}=U(x^p)$. The first equality holds for polynomials by
repeated use of $(a+b)^p=a^p+b^p$; it holds for power series because each coefficient of $U(x)^p$ involves only finitely
many coefficients of $U$.
\end{proof}

For $p=2$ this says that squaring a series modulo $2$ doubles every exponent: for instance,
$(x+x^3)^2\equiv x^2+x^6\pmod2$. In particular, the square of any series has, modulo $2$, no terms of odd degree.

\subsection{The Catalan numbers modulo 2}
\begin{lemma}\label{lem:catalan}
In every commutative ring the equation $U=x+U^2$ has exactly one solution with $U(0)=0$. Over the integers it is the
Catalan series $C(x)=x+x^2+2x^3+5x^4+14x^5+\cdots$, and
\[
C(x)\equiv x+x^2+x^4+x^8+\cdots=\sum_{k\ge0}x^{2^k}\pmod 2 .
\]
So $\coeff{n}{C}$ is odd exactly when $n$ is a power of $2$.
\end{lemma}

\begin{proof}
The coefficient of $x^n$ in $U^2$ involves only $u_1,\dots,u_{n-1}$, so the equation determines $u_n$ one at a time:
$u_1=1$ and $u_n=\sum_{i=1}^{n-1}u_iu_{n-i}$. Over $\F_2$, the series $S=\sum_kx^{2^k}$ satisfies
$x+S^2=x+\sum_kx^{2^{k+1}}=S$ by Lemma~\ref{lem:frobenius}, so by uniqueness $C\equiv S\pmod2$.
\end{proof}
The parity of the Catalan numbers is classical; see for example~\cite{deutschsagan}. (Here the Catalan number
$\frac1{n}\binom{2n-2}{n-1}$ appears as the coefficient of $x^n$, so ``$n$ a power of $2$'' corresponds to the familiar
statement that the $k$-th Catalan number is odd exactly when $k+1$ is a power of $2$.)

\section{Results}\label{sec:results}

Throughout, $C$ is the Catalan series, the unique normalized series with $C=x+C^2$ (Lemma~\ref{lem:catalan}), and
$c_n=\coeff{n}{C}$. The defining equation of the sequence is
\begin{equation}\label{eq:def}
A(x-x^2)=x+A(x)^2 .
\end{equation}

\begin{theorem}\label{thm:main}
\begin{enumerate}[label=(\alph*)]
\item Equation~\eqref{eq:def} has exactly one normalized solution $A$ with integer coefficients. (More generally, it
has exactly one normalized solution with coefficients in any commutative ring.) The sequence A393172 is
$a(n)=\coeff{n}{A}$.
\item Modulo $4$,
\[
A(x)\equiv 2C(x)-x\pmod 4 .
\]
\end{enumerate}
\end{theorem}

\begin{corollary}[Hanna's conjecture]\label{cor:conj}
For every $n>1$: if $n$ is a power of $2$ then $a(n)\equiv2\pmod4$, and otherwise $a(n)$ is divisible by $4$.
\end{corollary}

\begin{proof}
For $n>1$ the coefficient of $x^n$ in $2C(x)-x$ is $2c_n$, so Theorem~\ref{thm:main}(b) gives $a(n)\equiv2c_n\pmod4$.
If $c_n$ is odd, then $2c_n\equiv2\pmod4$; if $c_n$ is even, then $2c_n\equiv0\pmod4$. By Lemma~\ref{lem:catalan},
$c_n$ is odd exactly when $n$ is a power of $2$.
\end{proof}

For example, $c_8=429$ is odd, and indeed $a(8)=112290=4\cdot28072+2$; while $c_5=14$ is even, and $a(5)=260=4\cdot65$.

\section{Proofs}\label{sec:proofs}

\para{The idea in brief.} The substitution $x\mapsto x-x^2$ is undone by the Catalan series, because
$C(x)-C(x)^2=x$. Substituting $C$ into \eqref{eq:def} therefore turns it into the equation $A=C+A(C)^2$, which fixes
the coefficients of $A$ one at a time; so there is exactly one solution, in every commutative ring. For the congruence we
take the candidate $D=2C-x$ and check that it satisfies \eqref{eq:def} modulo $4$. Two short computations do this:
substituting $x-x^2$ into $D$ gives exactly $x+x^2$, and squaring $D$ gives $x^2$ plus multiples of $4$. By uniqueness,
$A$ and $D$ agree modulo $4$.

\begin{lemma}[An equation with a square]\label{lem:square}
Let $T$ and $P$ be normalized series with coefficients in a commutative ring $R$. There is exactly one normalized
series $G$ with coefficients in $R$ such that
\[
G(x)=T(x)+\bigl(G(P(x))\bigr)^2 .
\]
If $T$ and $P$ have integer coefficients, then so does $G$, and reducing $G$ modulo $m$ gives the solution of the same
equation over $\Z/m\Z$ (with $T$ and $P$ reduced modulo $m$).
\end{lemma}

\begin{proof}
The point is that the coefficient of $x^n$ on the right-hand side depends only on the coefficients of $G$ below degree
$n$. To see this, let $G$ and $H$ be two series with constant term $0$ that agree below degree $r$, so $G-H=O(x^r)$.
Since $P$ starts with $x$, the power $P^k$ has no terms below $x^k$, and so $G(P)-H(P)=O(x^r)$ as well. Moreover
$G(P)+H(P)$ has constant term $0$. Hence
\[
G(P)^2-H(P)^2=\bigl(G(P)-H(P)\bigr)\bigl(G(P)+H(P)\bigr)=O(x^{r+1}).
\]
In other words, changing $G$ in degree $r$ and above changes the right-hand side only in degree $r+1$ and above.

So the equation fixes the coefficients of $G$ one at a time. The constant term is $0$, the coefficient of $x$ is the
coefficient of $x$ in $T$, which is $1$, and for $n\ge2$ the coefficient of $x^n$ in $G$ must equal the coefficient of
$x^n$ in $T+G(P)^2$, which has already been computed from the earlier coefficients. Only additions and multiplications
are used, so this works in every commutative ring, and gives integer coefficients when $T$ and $P$ have them. Finally,
reducing the integer solution modulo $m$ gives a normalized series that satisfies the reduced equation (reduction is
compatible with all the operations involved), and by uniqueness it is \emph{the} solution over $\Z/m\Z$.
\end{proof}

\begin{lemma}[Two forms of the equation]\label{lem:forms}
Let $R$ be a commutative ring and $A$ a normalized series over $R$. Then $A$ satisfies \eqref{eq:def} if and only if
\begin{equation}\label{eq:catalanform}
A(x)=C(x)+\bigl(A(C(x))\bigr)^2 .
\end{equation}
\end{lemma}

\begin{proof}
Write $N(x)=x-x^2$. Then $N(C(x))=C-C^2=x$. By Lemma~\ref{lem:inverse}, $C$ is the compositional inverse of $N$, so also
$C(N(x))=x$. We use two facts about substitution: it can be done in steps, $(U\circ V)\circ W=U\circ(V\circ W)$;
and substituting into a sum or a product is the same as adding or multiplying the substituted terms.

Suppose $A(N)=x+A^2$. Substitute $C$ for $x$ on both sides. The left side becomes $A(N(C))=A(x)$, and the right side
becomes $C+A(C)^2$. This is \eqref{eq:catalanform}. Conversely, if \eqref{eq:catalanform} holds, substitute $N$ for $x$:
the left side becomes $A(N)$, and the right side becomes $C(N)+A(C(N))^2=x+A^2$. This is \eqref{eq:def}.
\end{proof}

\begin{proof}[Proof of Theorem~\ref{thm:main}]
(a) By Lemma~\ref{lem:forms}, the normalized solutions of \eqref{eq:def} are exactly the normalized solutions of
\eqref{eq:catalanform}. By Lemma~\ref{lem:square} with $T=P=C$, there is exactly one, in every commutative ring. Over the
integers it has integer coefficients, because $C$ does.

(b) By part (a) applied to the ring $\Z/4\Z$, and the last sentence of Lemma~\ref{lem:square}, the reduction of $A$ modulo
$4$ is the only normalized series over $\Z/4\Z$ that satisfies \eqref{eq:def}. So it is enough to check that the normalized
series $D=2C-x$ satisfies \eqref{eq:def} in $\Z/4\Z$. On the left side, using $C(x-x^2)=x$,
\[
D(x-x^2)=2\,C(x-x^2)-(x-x^2)=2x-x+x^2=x+x^2 .
\]
On the right side,
\[
x+D^2=x+(2C-x)^2=x+4C^2-4xC+x^2=x+x^2\quad\text{in }\Z/4\Z,
\]
because the terms $4C^2$ and $4xC$ vanish modulo $4$. The two sides agree, so $A\equiv D=2C-x\pmod4$.
\end{proof}

\section{Checking the results}\label{sec:checks}

A mathematical proof can contain mistakes, so the results were checked in two independent ways.

\para{Numerical check.}
We solved \eqref{eq:def} coefficient by coefficient, with exact integer arithmetic. This reproduces all 24 terms
printed in the OEIS entry. We then checked that every coefficient of $A(x)-\bigl(2C(x)-x\bigr)$ below degree $200$ is
divisible by $4$, and that the rule of Corollary~\ref{cor:conj} holds for every $n$ with $1<n<200$.
The computation was done with the PARI/GP computer algebra system, using code that was written separately from the
formal proofs (file \texttt{verification/independent\_checks.gp} in the accompanying repository~\cite{repo}).
A computation can only test finitely many cases, so it is a check on the definitions and on the indexing, not a proof.

\para{Formal proof in Lean.}
Lean~\cite{lean} is a \emph{proof assistant}: a program in which mathematical statements and their proofs are written
in a precise formal language, and every step is checked mechanically down to the basic rules of logic. The theorems of Section~\ref{sec:results} are proved in Lean~4 (version 4.35.0) on top of
Mathlib~\cite{mathlib}, Lean's library of formalized mathematics. The statements are collected in the file \texttt{lean/Entries/A393172.lean} of the repository,
in the namespace \texttt{OEIS.A393172}. There the Catalan series is called \texttt{catalan}, and the main proofs are in the module \texttt{lean/Families/CatalanComposition.lean}.
The table lists where each result of this paper is proved.
\begin{center}\small\renewcommand{\arraystretch}{1.25}
\begin{tabular}{>{\raggedright\arraybackslash}p{0.52\linewidth}>{\raggedright\arraybackslash}p{0.4\linewidth}}
\toprule
Statement in this paper & Lean theorem(s)\\
\midrule
Theorem~\ref{thm:main}(a) for integer coefficients; $A$ is normalized and satisfies \eqref{eq:def} & \texttt{first\_\allowbreak{}existence}, \texttt{first\_\allowbreak{}definition} \\
Theorem~\ref{thm:main}(b): $A\equiv2C-x\pmod4$ & \texttt{complete\_\allowbreak{}series\_\allowbreak{}mod\_\allowbreak{}four} \\
Corollary~\ref{cor:conj}, also for any normalized integer solution of \eqref{eq:def} & \texttt{first\_\allowbreak{}residue}, \texttt{divisibility\_\allowbreak{}complement}, \texttt{first\_\allowbreak{}arbitrary\_\allowbreak{}source} \\
\bottomrule
\end{tabular}
\end{center}
The Lean project was rebuilt from source, and the axioms on which each of these theorems depends were listed with
Lean's \texttt{\#print axioms} command. Every one depends only on the three standard axioms of Lean's logic
(\texttt{propext}, \texttt{Classical.choice}, \texttt{Quot.sound}); no step is left unproved. A formal proof certifies
exactly the statement that Lean checks. Whether that statement says what the OEIS entry means is a separate question,
so the Lean definitions were also compared with the definition in the OEIS entry
(Section~\ref{sec:sequence}). The principal theorems are stated once more, using only Mathlib's definitions, in the
file \texttt{lean/Challenges/Catalan.lean}, and Lean's Comparator tool confirms that the
proved theorems have exactly these statements.

\section{Discussion}\label{sec:discussion}

\para{Scope.} Theorem~\ref{thm:main} determines the remainder of every coefficient $a(n)$ modulo $4$, so it settles
the conjecture of the entry completely. It says nothing about the size of $a(n)$, and nothing about remainders modulo
higher powers of $2$.

\para{Modulo 8.} A computer calculation suggests a similar, but more intricate, rule modulo $8$. For every $n$ with
$1<n<200$ we found: $a(n)\equiv2\pmod8$ when $n$ is a power of $2$; $a(n)\equiv4\pmod8$ when $n=2^k\,(4^j+1)$ for some
$j\ge1$ and $k\ge0$ (that is, $n=5,10,17,20,34,40,65,68,80,\dots$); and $a(n)\equiv0\pmod8$ otherwise. We have not tried
to prove this. A proof would probably need finer information about the Catalan numbers, namely their remainders
modulo $4$.

\para{The companion sequence.} The entry A393171~\cite{oeis171} is about the series $B=x+A^2$. Squaring is useful here:
if $A=D+2^kK$ for some series $K$ with integer coefficients and some $k\ge1$, then $A^2=D^2+2^{k+1}DK+2^{2k}K^2$, so $A^2$
is known modulo $2^{k+1}$ once $A$ is known modulo $2^k$. In this way Theorem~\ref{thm:main} determines $B$ modulo $8$,
which proves Hanna's conjecture for A393171 (see the companion paper). The rule modulo $8$ suggested above would, in the
same way, determine $B$ modulo $16$.

\section*{About this paper}\label{sec:about}
\para{How the work was done.}
This paper is one of 29 papers on conjectures from the OEIS that came out of an AI research effort directed by
Dev~D.~Goyal~\cite{repo}. The roles were as follows.
\begin{itemize}
\item \emph{Dev~D.~Goyal} started the project, chose which problems to pursue, asked for proofs, formal verification and
revisions, and decided what to publish. Dev~D.~Goyal did not supply proof strategies or intermediate arguments.
\item \emph{ChatGPT} (OpenAI), using the GPT-6 model with Pro reasoning, worked autonomously from those prompts. It
searched the OEIS for suitable conjectures, found the proofs, wrote the Lean~4 formalization, and wrote a first draft
of this paper.
\item \emph{Claude Opus~5.5}, used with Claude Code, checked the work independently. It rebuilt the Lean project from source and confirmed that the theorems listed in Section~\ref{sec:checks} rely only on
Lean's standard axioms; compared the Lean definitions with the OEIS definition; recomputed the sequence with separately
written PARI/GP code, compared it with the terms published in the OEIS, and tested every claim of the paper on many
terms; and read every step of the proofs. It then wrote the present version of the paper, expanding the
explanations and, in places, writing out an argument in more detail or by a more elementary route. The theorems of
Section~\ref{sec:results} and the proof strategies are ChatGPT's. The Discussion section, including the short arguments
in it, was written by Claude Opus~5.5; its observations were checked by computer and are not part of the Lean
formalization.
\end{itemize}
Dev~D.~Goyal has reviewed the proofs and is responsible for this paper; this review is not an independent expert
validation.

\para{Review status.} This paper has not been peer reviewed, and it has not been submitted to a journal or to the OEIS.

\para{Priority.} An OEIS entry that labels a statement as a conjecture shows only that its contributor did not know
a proof; it does not show that no proof existed. During the research only limited searches for earlier work were
made. No claim of priority is made.

\para{Corrections.} Please report any error in the issue tracker of the repository~\cite{repo}, where the Lean sources and the numerical checks are also
available.

\section*{Acknowledgements}
The sequence, its formulas and the conjecture are due to Paul~D.~Hanna~\cite{oeis}. The recurrence in the entry was
contributed by Seiichi~Manyama.

\begin{thebibliography}{99}\small
\raggedright
\bibitem{oeis} P.~D.~Hanna, entry A393172 (G.f.\ $A(x)$ satisfies $A(x-x^2)=x+A(x)^2$), \emph{The On-Line Encyclopedia of Integer Sequences}, February 2026. \url{https://oeis.org/A393172}. Text as of October 2026.
\bibitem{oeis171} P.~D.~Hanna, entry A393171 (G.f.\ satisfies $A(x)=x+A(C(x))^2$, where $C(x)=x+C(x)^2$), \emph{The On-Line Encyclopedia of Integer Sequences}, February 2026. \url{https://oeis.org/A393171}. Text as of October 2026.
\bibitem{deutschsagan} E.~Deutsch and B.~E.~Sagan, Congruences for Catalan and Motzkin numbers and related sequences, \emph{Journal of Number Theory} 117 (2006), 191--215.
\bibitem{oeisdb} N.~J.~A.~Sloane and The OEIS Foundation Inc., \emph{The On-Line Encyclopedia of Integer Sequences}. \url{https://oeis.org}.
\bibitem{stanley} R.~P.~Stanley, \emph{Enumerative Combinatorics}, Volume~2, Cambridge University Press, 1999. (Chapter~5 treats formal power series, composition and compositional inverses.)
\bibitem{niven} I.~Niven, Formal power series, \emph{American Mathematical Monthly} 76 (1969), 871--889.
\bibitem{lean} L.~de~Moura and S.~Ullrich, The Lean~4 theorem prover and programming language, in \emph{Automated Deduction -- CADE 28}, Lecture Notes in Computer Science 12699, Springer, 2021, 625--635.
\bibitem{mathlib} The mathlib Community, The Lean mathematical library, in \emph{Proceedings of the 9th ACM SIGPLAN International Conference on Certified Programs and Proofs (CPP 2020)}, 367--381. Version used: commit \texttt{56a005855ea1}.
\bibitem{repo} D.~D.~Goyal, \emph{Proofs and corrections of OEIS conjectures}, repository, 2026. \url{https://github.com/dvgyl/oeis-proofs}.
\end{thebibliography}
\end{document}
